\begin{abstract}
Version control works for source code because its merge understands the medium:
lines. Knowledge graphs shared between people and software agents are not lines,
and merging their serialisations makes conflicts out of ordering and misses the
ones that matter. We describe and evaluate a three-way merge for RDF in which
the \emph{schema} decides what a conflict is: the merge is set algebra over
canonical triples, and a slot is contended only where a SHACL shape declares
that it can hold at most one value. Multi-valued predicates union; functional
predicates with divergent values are handed to a person. The same shapes graph
then validates the merge result, so the schema that defines conflicts also
audits their resolution.

We do not claim the first three-way merge for RDF. Git-backed RDF versioning
with merge strategies is established prior work, and the closest neighbour ---
Quit Store's Context Merge --- is included here as a baseline that this operator
extends rather than replaces. Our contribution is an implementation in a
production multi-agent store and an evaluation of what schema-derived conflict
semantics buys against that neighbour and against six other strategies.

On a synthetic divergence benchmark over five seeds, the shape-aware operator
raised $33$ conflicts, all true positives, with no false positives, no triples
lost, none fabricated, and no SHACL violations admitted. The context-overlap
baseline detected the \emph{same} $33$ conflicts and asked $845$ questions to do
it. We also report a limit no schema removes: when two sides mint different
names for one entity, the divergence is invisible to any triple-level operator,
including this one. Replaying a recorded production corpus against the shipped
operator --- $105$ repairs a person actually performed, $93$ of them after
excluding $12$ chained pairs, and $939$ real duplicate-value incidents ---
reproduces both the blindness ($0$ of $93$ alias repairs, matching $0$ of $21$
on the synthetic arm) and the benefit, on data nobody generated for the
purpose.
\end{abstract}
